% Source: Confluence page 12025880, version 6; FR/NFR content retained.
\chapter{Functional and Non-Functional Requirements}
\label{ch:requirements}

This chapter defines Trippa's functional (FR) and non-functional (NFR) requirements, retaining their source identifiers for architectural reference. Core planning remains independent of optional AI assistance, and mobile clients use the same backend rules and trip data. Unapproved thresholds and unfinished system models remain marked TBD.

\section{Functional Requirements}
The eleven functional areas include mobile planning and management. They specify intended capabilities, not completed implementation.

\Needspace*{12\baselineskip}
\subsection{FR-1: User Authentication and Account Management}

The system shall support guest access, authenticated user accounts, and role-based authorization.


\begingroup
\small
\setlength{\tabcolsep}{5pt}
\renewcommand{\arraystretch}{1.18}
\begin{longtable}{@{}P{\dimexpr 0.16\linewidth-0.32\tabcolsep\relax}P{\dimexpr 0.84\linewidth-1.68\tabcolsep\relax}@{}}
\caption{User Authentication and Account Management}\label{tab:req-3}\\
\toprule
\textbf{ID} & \textbf{Functional Requirement} \\
\midrule
\endfirsthead
\multicolumn{2}{@{}l}{\small\itshape Table \thetable{} (continued)}\\
\toprule
\textbf{ID} & \textbf{Functional Requirement} \\
\midrule
\endhead
\midrule
\multicolumn{2}{r@{}}{\footnotesize\itshape Continued on next page}\\
\endfoot
\bottomrule
\endlastfoot
FR-1.1 & The system shall allow guests to access the basic trip-planning workflow without authentication. \\ \addlinespace[2pt]
FR-1.2 & The system shall allow guests to register an account using supported registration methods. \\ \addlinespace[2pt]
FR-1.3 & The system shall allow registered users to log in and log out. \\ \addlinespace[2pt]
FR-1.4 & The system shall authenticate users before granting access to protected account-specific functions. \\ \addlinespace[2pt]
FR-1.5 & The system shall allow registered users to view and update their profile information. \\ \addlinespace[2pt]
FR-1.6 & The system shall provide an account recovery mechanism. \\ \addlinespace[2pt]
FR-1.7 & The system shall maintain authenticated sessions according to the supported authentication mechanism. \\ \addlinespace[2pt]
FR-1.8 & The system shall distinguish between regular users and administrators and enforce the appropriate access permissions. \\ \addlinespace[2pt]
FR-1.9 & The system shall prevent users from accessing or modifying another user's private account information and saved trips without authorization. \\ \addlinespace[2pt]
\end{longtable}
\endgroup


\Needspace*{12\baselineskip}
\subsection{FR-2: User Preferences Management}

The system shall support reusable planning preferences to reduce repetitive input when creating new trips.


\begingroup
\small
\setlength{\tabcolsep}{5pt}
\renewcommand{\arraystretch}{1.18}
\begin{longtable}{@{}P{\dimexpr 0.16\linewidth-0.32\tabcolsep\relax}P{\dimexpr 0.84\linewidth-1.68\tabcolsep\relax}@{}}
\caption{User Preferences Management}\label{tab:req-4}\\
\toprule
\textbf{ID} & \textbf{Functional Requirement} \\
\midrule
\endfirsthead
\multicolumn{2}{@{}l}{\small\itshape Table \thetable{} (continued)}\\
\toprule
\textbf{ID} & \textbf{Functional Requirement} \\
\midrule
\endhead
\midrule
\multicolumn{2}{r@{}}{\footnotesize\itshape Continued on next page}\\
\endfoot
\bottomrule
\endlastfoot
FR-2.1 & The system shall allow registered users to configure their default transportation mode. \\ \addlinespace[2pt]
FR-2.2 & The system shall allow registered users to configure their preferred spending range or default budget. \\ \addlinespace[2pt]
FR-2.3 & The system shall support storing additional planning preferences, including activity interests and preferred trip pace, when enabled. \\ \addlinespace[2pt]
FR-2.4 & The system shall persist user preferences and retrieve them in subsequent authenticated sessions. \\ \addlinespace[2pt]
FR-2.5 & The system shall automatically apply available saved preferences when creating a new trip request. \\ \addlinespace[2pt]
FR-2.6 & The system shall allow users to override default preferences for an individual trip without changing their saved preferences. \\ \addlinespace[2pt]
FR-2.7 & The system shall allow registered users to update their saved preferences. \\ \addlinespace[2pt]
\end{longtable}
\endgroup


\Needspace*{12\baselineskip}
\subsection{FR-3: Trip Request Creation and Management}

The system shall allow users to define, review, and modify the constraints required for itinerary generation.


\begingroup
\small
\setlength{\tabcolsep}{5pt}
\renewcommand{\arraystretch}{1.18}
\begin{longtable}{@{}P{\dimexpr 0.16\linewidth-0.32\tabcolsep\relax}P{\dimexpr 0.84\linewidth-1.68\tabcolsep\relax}@{}}
\caption{Trip Request Creation and Management}\label{tab:req-5}\\
\toprule
\textbf{ID} & \textbf{Functional Requirement} \\
\midrule
\endfirsthead
\multicolumn{2}{@{}l}{\small\itshape Table \thetable{} (continued)}\\
\toprule
\textbf{ID} & \textbf{Functional Requirement} \\
\midrule
\endhead
\midrule
\multicolumn{2}{r@{}}{\footnotesize\itshape Continued on next page}\\
\endfoot
\bottomrule
\endlastfoot
FR-3.1 & The system shall allow guests and registered users to create a trip request. \\ \addlinespace[2pt]
FR-3.2 & The system shall allow users to specify the trip date, starting time, and ending time. \\ \addlinespace[2pt]
FR-3.3 & The system shall allow users to select a starting location and an optional ending location. \\ \addlinespace[2pt]
FR-3.4 & The system shall allow users to specify whether the trip should return to its starting location. \\ \addlinespace[2pt]
FR-3.5 & The system shall allow users to select the intended transportation mode. \\ \addlinespace[2pt]
FR-3.6 & The system shall allow users to specify an available budget for the trip. \\ \addlinespace[2pt]
FR-3.7 & The system shall allow users to specify preferred activities, interests, or desired experiences. \\ \addlinespace[2pt]
FR-3.8 & The system shall allow users to select desired places and assign priorities or required-place constraints where supported. \\ \addlinespace[2pt]
FR-3.9 & The system shall allow users to review and modify the complete trip request before generating an itinerary. \\ \addlinespace[2pt]
FR-3.10 & The system shall validate required fields and reject invalid or inconsistent trip request data with understandable feedback. \\ \addlinespace[2pt]
FR-3.11 & The system shall allow users to revise a trip request and regenerate the corresponding itinerary. \\ \addlinespace[2pt]
\end{longtable}
\endgroup


\Needspace*{12\baselineskip}
\subsection{FR-4: Place Discovery and Map Interaction}

The system shall integrate with external geographic services to support location discovery, route visualization, and map-based itinerary interaction.


\begingroup
\small
\setlength{\tabcolsep}{5pt}
\renewcommand{\arraystretch}{1.18}
\begin{longtable}{@{}P{\dimexpr 0.16\linewidth-0.32\tabcolsep\relax}P{\dimexpr 0.84\linewidth-1.68\tabcolsep\relax}@{}}
\caption{Place Discovery and Map Interaction}\label{tab:req-6}\\
\toprule
\textbf{ID} & \textbf{Functional Requirement} \\
\midrule
\endfirsthead
\multicolumn{2}{@{}l}{\small\itshape Table \thetable{} (continued)}\\
\toprule
\textbf{ID} & \textbf{Functional Requirement} \\
\midrule
\endhead
\midrule
\multicolumn{2}{r@{}}{\footnotesize\itshape Continued on next page}\\
\endfoot
\bottomrule
\endlastfoot
FR-4.1 & The system shall allow users to search for places and geographic locations. \\ \addlinespace[2pt]
FR-4.2 & The system shall allow users to select search results as starting locations, ending locations, or desired places. \\ \addlinespace[2pt]
FR-4.3 & The system shall display selected locations and scheduled activities on an interactive map. \\ \addlinespace[2pt]
FR-4.4 & The system shall display available place information relevant to itinerary planning. \\ \addlinespace[2pt]
FR-4.5 & The system shall retrieve route information, estimated travel durations, and distances between itinerary locations. \\ \addlinespace[2pt]
FR-4.6 & The system shall display the planned route between consecutive activities. \\ \addlinespace[2pt]
FR-4.7 & The system shall synchronize activity selection between the map and the itinerary summary. \\ \addlinespace[2pt]
FR-4.8 & The system shall allow users to access supported itinerary editing actions through map or itinerary interactions. \\ \addlinespace[2pt]
FR-4.9 & The system shall support opening selected destinations or routes in an external navigation service. \\ \addlinespace[2pt]
\end{longtable}
\endgroup


\Needspace*{12\baselineskip}
\subsection{FR-5: Automatic Itinerary Generation and Feasibility Validation}

The system shall generate structured itineraries using user constraints and verify their feasibility before they are accepted.


\begingroup
\small
\setlength{\tabcolsep}{5pt}
\renewcommand{\arraystretch}{1.18}
\begin{longtable}{@{}P{\dimexpr 0.16\linewidth-0.32\tabcolsep\relax}P{\dimexpr 0.84\linewidth-1.68\tabcolsep\relax}@{}}
\caption{Automatic Itinerary Generation and Feasibility Validation}\label{tab:req-7}\\
\toprule
\textbf{ID} & \textbf{Functional Requirement} \\
\midrule
\endfirsthead
\multicolumn{2}{@{}l}{\small\itshape Table \thetable{} (continued)}\\
\toprule
\textbf{ID} & \textbf{Functional Requirement} \\
\midrule
\endhead
\midrule
\multicolumn{2}{r@{}}{\footnotesize\itshape Continued on next page}\\
\endfoot
\bottomrule
\endlastfoot
FR-5.1 & The system shall generate an ordered candidate itinerary from a valid trip request. \\ \addlinespace[2pt]
FR-5.2 & The system shall identify candidate places based on the user's locations, preferences, and trip constraints. \\ \addlinespace[2pt]
FR-5.3 & The system shall estimate the duration and expected cost of each proposed activity using available information. \\ \addlinespace[2pt]
FR-5.4 & The system shall calculate the estimated travel time between consecutive activities. \\ \addlinespace[2pt]
FR-5.5 & The system shall assign planned arrival, start, and end times to itinerary activities. \\ \addlinespace[2pt]
FR-5.6 & The system shall consider the selected transportation mode when estimating travel times and constructing the itinerary. \\ \addlinespace[2pt]
FR-5.7 & The system shall verify that the itinerary fits within the specified trip time window. \\ \addlinespace[2pt]
FR-5.8 & The system shall verify compatibility with known place opening hours. \\ \addlinespace[2pt]
FR-5.9 & The system shall evaluate the itinerary against the user's specified budget constraints and cost-estimation policy. \\ \addlinespace[2pt]
FR-5.10 & The system shall verify that required places are handled according to their assigned priority and inclusion rules. \\ \addlinespace[2pt]
FR-5.11 & The system shall reject or identify candidate itineraries that violate mandatory feasibility constraints. \\ \addlinespace[2pt]
FR-5.12 & The system shall provide understandable feedback when no feasible itinerary can be generated. \\ \addlinespace[2pt]
FR-5.13 & The system shall display the generated itinerary through an interactive map and structured Trip Summary. \\ \addlinespace[2pt]
FR-5.14 & The system shall present estimated travel durations, activity durations, and expected costs to the user. \\ \addlinespace[2pt]
\end{longtable}
\endgroup


\Needspace*{12\baselineskip}
\subsection{FR-6: Itinerary Editing and Validation}

The system shall allow users to modify a generated itinerary while ensuring that the resulting schedule remains consistent with the trip constraints.


\begingroup
\small
\setlength{\tabcolsep}{5pt}
\renewcommand{\arraystretch}{1.18}
\begin{longtable}{@{}P{\dimexpr 0.16\linewidth-0.32\tabcolsep\relax}P{\dimexpr 0.84\linewidth-1.68\tabcolsep\relax}@{}}
\caption{Itinerary Editing and Validation}\label{tab:req-8}\\
\toprule
\textbf{ID} & \textbf{Functional Requirement} \\
\midrule
\endfirsthead
\multicolumn{2}{@{}l}{\small\itshape Table \thetable{} (continued)}\\
\toprule
\textbf{ID} & \textbf{Functional Requirement} \\
\midrule
\endhead
\midrule
\multicolumn{2}{r@{}}{\footnotesize\itshape Continued on next page}\\
\endfoot
\bottomrule
\endlastfoot
FR-6.1 & The system shall allow users to add a new visit to an existing itinerary. \\ \addlinespace[2pt]
FR-6.2 & The system shall allow users to remove a planned visit. \\ \addlinespace[2pt]
FR-6.3 & The system shall allow users to modify supported visit properties, including planned time, duration, and estimated budget. \\ \addlinespace[2pt]
FR-6.4 & The system shall support changing the order of itinerary activities. \\ \addlinespace[2pt]
FR-6.5 & The system shall recalculate affected travel segments and activity schedules after itinerary modifications. \\ \addlinespace[2pt]
FR-6.6 & The system shall revalidate modified itineraries against applicable feasibility constraints. \\ \addlinespace[2pt]
FR-6.7 & The system shall display validation errors or conflicts caused by proposed modifications. \\ \addlinespace[2pt]
FR-6.8 & The system shall allow users to review the updated itinerary before accepting it. \\ \addlinespace[2pt]
FR-6.9 & The system shall allow users to accept a feasible itinerary as the plan to be followed. \\ \addlinespace[2pt]
\end{longtable}
\endgroup


\Needspace*{12\baselineskip}
\subsection{FR-7: Temporary Storage, Saved Trips, and Trip History}

The system shall support temporary guest planning and persistent trip management for registered users.


\begingroup
\small
\setlength{\tabcolsep}{5pt}
\renewcommand{\arraystretch}{1.18}
\begin{longtable}{@{}P{\dimexpr 0.16\linewidth-0.32\tabcolsep\relax}P{\dimexpr 0.84\linewidth-1.68\tabcolsep\relax}@{}}
\caption{Temporary Storage, Saved Trips, and Trip History}\label{tab:req-9}\\
\toprule
\textbf{ID} & \textbf{Functional Requirement} \\
\midrule
\endfirsthead
\multicolumn{2}{@{}l}{\small\itshape Table \thetable{} (continued)}\\
\toprule
\textbf{ID} & \textbf{Functional Requirement} \\
\midrule
\endhead
\midrule
\multicolumn{2}{r@{}}{\footnotesize\itshape Continued on next page}\\
\endfoot
\bottomrule
\endlastfoot
FR-7.1 & The system shall allow guests to temporarily retain generated itineraries without authentication. \\ \addlinespace[2pt]
FR-7.2 & The system shall allow guests to view and modify temporarily retained itineraries while the temporary data remains available. \\ \addlinespace[2pt]
FR-7.3 & The system shall require authentication before permanently saving an itinerary to a user account. \\ \addlinespace[2pt]
FR-7.4 & The system shall allow guests to initiate login or registration when requesting to save an itinerary. \\ \addlinespace[2pt]
FR-7.5 & The system shall preserve the current temporary itinerary through the supported authentication flow so that it can be saved afterward. \\ \addlinespace[2pt]
FR-7.6 & The system shall allow authenticated users to save, retrieve, and manage their itineraries. \\ \addlinespace[2pt]
FR-7.7 & The system shall maintain trip lifecycle states, including Ready, Ongoing, and Completed. \\ \addlinespace[2pt]
FR-7.8 & The system shall allow registered users to view saved trips with their relevant status, date, and summary information. \\ \addlinespace[2pt]
FR-7.9 & The system shall direct users to the appropriate trip interface according to the selected trip's lifecycle state. \\ \addlinespace[2pt]
FR-7.10 & The system shall preserve completed trip records for future reference. \\ \addlinespace[2pt]
FR-7.11 & The system shall allow registered users to delete their saved trips. \\ \addlinespace[2pt]
\end{longtable}
\endgroup


\Needspace*{12\baselineskip}
\subsection{FR-8: Active Trip Execution and Dynamic Replanning}

The system shall support trip execution and adjust the unfinished itinerary when actual progress differs from the original schedule.


\begingroup
\small
\setlength{\tabcolsep}{5pt}
\renewcommand{\arraystretch}{1.18}
\begin{longtable}{@{}P{\dimexpr 0.16\linewidth-0.32\tabcolsep\relax}P{\dimexpr 0.84\linewidth-1.68\tabcolsep\relax}@{}}
\caption{Active Trip Execution and Dynamic Replanning}\label{tab:req-10}\\
\toprule
\textbf{ID} & \textbf{Functional Requirement} \\
\midrule
\endfirsthead
\multicolumn{2}{@{}l}{\small\itshape Table \thetable{} (continued)}\\
\toprule
\textbf{ID} & \textbf{Functional Requirement} \\
\midrule
\endhead
\midrule
\multicolumn{2}{r@{}}{\footnotesize\itshape Continued on next page}\\
\endfoot
\bottomrule
\endlastfoot
FR-8.1 & The system shall allow users to start an accepted trip and transition it to the Ongoing state. \\ \addlinespace[2pt]
FR-8.2 & The system shall display the current visit, completed visits, and remaining scheduled activities. \\ \addlinespace[2pt]
FR-8.3 & The system shall allow users to mark visits as completed. \\ \addlinespace[2pt]
FR-8.4 & The system shall allow users to skip planned visits. \\ \addlinespace[2pt]
FR-8.5 & The system shall allow users to report changes affecting unfinished activities, including delays or changes in visit duration. \\ \addlinespace[2pt]
FR-8.6 & The system shall record trip progress events and update the corresponding trip state. \\ \addlinespace[2pt]
FR-8.7 & The system shall allow users to modify supported unfinished activities during an ongoing trip. \\ \addlinespace[2pt]
FR-8.8 & The system shall reassess the remaining itinerary when a reported event or modification affects its feasibility. \\ \addlinespace[2pt]
FR-8.9 & The system shall generate an adjusted itinerary for the unfinished portion of the trip when replanning is necessary. \\ \addlinespace[2pt]
FR-8.10 & The system shall preserve completed activities and their recorded history during replanning. \\ \addlinespace[2pt]
FR-8.11 & The system shall reuse applicable scheduling and feasibility rules when generating an adjusted itinerary. \\ \addlinespace[2pt]
FR-8.12 & The system shall present proposed changes to the remaining itinerary for review when confirmation is required. \\ \addlinespace[2pt]
FR-8.13 & The system shall update the active itinerary after the proposed changes are accepted. \\ \addlinespace[2pt]
FR-8.14 & The system shall inform users when the remaining itinerary cannot satisfy mandatory constraints. \\ \addlinespace[2pt]
FR-8.15 & The system shall allow users to end an ongoing trip and record its completion. \\ \addlinespace[2pt]
\end{longtable}
\endgroup


\Needspace*{12\baselineskip}
\subsection{FR-9: Administration and System Management}

The system shall provide administrative capabilities for authorized personnel to maintain user accounts and monitor system operations.


\begingroup
\small
\setlength{\tabcolsep}{5pt}
\renewcommand{\arraystretch}{1.18}
\begin{longtable}{@{}P{\dimexpr 0.16\linewidth-0.32\tabcolsep\relax}P{\dimexpr 0.84\linewidth-1.68\tabcolsep\relax}@{}}
\caption{Administration and System Management}\label{tab:req-11}\\
\toprule
\textbf{ID} & \textbf{Functional Requirement} \\
\midrule
\endfirsthead
\multicolumn{2}{@{}l}{\small\itshape Table \thetable{} (continued)}\\
\toprule
\textbf{ID} & \textbf{Functional Requirement} \\
\midrule
\endhead
\midrule
\multicolumn{2}{r@{}}{\footnotesize\itshape Continued on next page}\\
\endfoot
\bottomrule
\endlastfoot
FR-9.1 & The system shall provide a protected administration interface. \\ \addlinespace[2pt]
FR-9.2 & The system shall allow administrators to view and search registered user accounts. \\ \addlinespace[2pt]
FR-9.3 & The system shall allow administrators to manage user account status. \\ \addlinespace[2pt]
FR-9.4 & The system shall allow administrators to view relevant operational information and system error records. \\ \addlinespace[2pt]
FR-9.5 & The system shall restrict administrative functions to authorized administrators. \\ \addlinespace[2pt]
FR-9.6 & The system shall prevent administrative access from bypassing established privacy and authorization rules without an explicitly defined permission. \\ \addlinespace[2pt]
\end{longtable}
\endgroup


\Needspace*{12\baselineskip}
\subsection{FR-10: AI-Assisted Trip Planning and Modification}

The system shall support AI-assisted interactions as an extension of the core planning functionality.


\begingroup
\small
\setlength{\tabcolsep}{5pt}
\renewcommand{\arraystretch}{1.18}
\begin{longtable}{@{}P{\dimexpr 0.16\linewidth-0.32\tabcolsep\relax}P{\dimexpr 0.84\linewidth-1.68\tabcolsep\relax}@{}}
\caption{AI-Assisted Trip Planning and Modification}\label{tab:req-12}\\
\toprule
\textbf{ID} & \textbf{Functional Requirement} \\
\midrule
\endfirsthead
\multicolumn{2}{@{}l}{\small\itshape Table \thetable{} (continued)}\\
\toprule
\textbf{ID} & \textbf{Functional Requirement} \\
\midrule
\endhead
\midrule
\multicolumn{2}{r@{}}{\footnotesize\itshape Continued on next page}\\
\endfoot
\bottomrule
\endlastfoot
FR-10.1 & The system shall allow users to submit supported trip-planning requests using natural language. \\ \addlinespace[2pt]
FR-10.2 & The system shall interpret natural-language requests and extract relevant planning information, such as time, budget, location, and preferences. \\ \addlinespace[2pt]
FR-10.3 & The system shall convert interpreted requests into structured data compatible with the core planning functionality. \\ \addlinespace[2pt]
FR-10.4 & The system shall allow users to request supported itinerary modifications using natural-language commands. \\ \addlinespace[2pt]
FR-10.5 & The system shall identify supported modification intentions, such as adding, removing, or reordering visits. \\ \addlinespace[2pt]
FR-10.6 & The system shall validate AI-generated structured suggestions before applying them to an itinerary. \\ \addlinespace[2pt]
FR-10.7 & The system shall use the core planning and feasibility functionality to evaluate proposed AI-assisted modifications. \\ \addlinespace[2pt]
FR-10.8 & The system shall provide understandable feedback when an AI request cannot be interpreted or applied. \\ \addlinespace[2pt]
FR-10.9 & The system shall preserve the availability of core non-AI trip-planning functions when the AI service is unavailable. \\ \addlinespace[2pt]
\end{longtable}
\endgroup



\Needspace*{12\baselineskip}
\subsection{FR-11: Mobile Trip Planning and Management}

\begingroup
\small
\setlength{\tabcolsep}{5pt}
\renewcommand{\arraystretch}{1.18}
\begin{longtable}{@{}P{\dimexpr 0.16\linewidth-0.32\tabcolsep\relax}P{\dimexpr 0.84\linewidth-1.68\tabcolsep\relax}@{}}
\caption{Mobile Trip Planning and Management}\label{tab:req-23}\\
\toprule
\textbf{ID} & \textbf{Functional Requirement} \\
\midrule
\endfirsthead
\multicolumn{2}{@{}l}{\small\itshape Table \thetable{} (continued)}\\
\toprule
\textbf{ID} & \textbf{Functional Requirement} \\
\midrule
\endhead
\midrule
\multicolumn{2}{r@{}}{\footnotesize\itshape Continued on next page}\\
\endfoot
\bottomrule
\endlastfoot
FR-11.1 & The system shall allow users to complete the trip creation workflow on a mobile device. \\ \addlinespace[2pt]
FR-11.2 & The system shall support touch-based selection of start locations, end locations, and desired places. \\ \addlinespace[2pt]
FR-11.3 & The system shall allow users to review and correct trip request information before generating an itinerary on mobile. \\ \addlinespace[2pt]
FR-11.4 & The system shall provide a map-first mobile itinerary view. \\ \addlinespace[2pt]
FR-11.5 & The system shall allow users to open and close the Trip Summary while preserving the current map context. \\ \addlinespace[2pt]
FR-11.6 & The system shall allow users to edit visit information, including time, duration, and budget, within the mobile Trip Summary. \\ \addlinespace[2pt]
FR-11.7 & The system shall allow users to add and remove visits through the mobile interface. \\ \addlinespace[2pt]
FR-11.8 & The system shall allow users to perform supported active-trip actions, including completing visits, skipping visits, requesting replanning, and ending trips. \\ \addlinespace[2pt]
FR-11.9 & The system shall allow users to review updated remaining itineraries before confirming applicable replanning changes. \\ \addlinespace[2pt]
FR-11.10 & The system shall provide mobile access to saved trips, ongoing trips, and completed trip history. \\ \addlinespace[2pt]
FR-11.11 & The system shall use the same backend validation and authoritative trip information for desktop and mobile operations. \\ \addlinespace[2pt]
\end{longtable}
\endgroup



\section{Non-Functional Requirements}


Non-functional requirements specify the quality attributes necessary to ensure that Trippa is usable, reliable, secure, maintainable, and capable of supporting its intended workload.


Where quantitative acceptance thresholds have not yet been approved, they are marked as TBD.


\Needspace*{12\baselineskip}
\subsection{NFR-1: Performance Efficiency}

\begingroup
\small
\setlength{\tabcolsep}{5pt}
\renewcommand{\arraystretch}{1.18}
\begin{longtable}{@{}P{\dimexpr 0.15\linewidth-0.60\tabcolsep\relax}P{\dimexpr 0.49\linewidth-1.96\tabcolsep\relax}P{\dimexpr 0.36\linewidth-1.44\tabcolsep\relax}@{}}
\caption{Performance Efficiency}\label{tab:req-13}\\
\toprule
\textbf{ID} & \textbf{Non-Functional Requirement} & \textbf{Acceptance Criteria} \\
\midrule
\endfirsthead
\multicolumn{3}{@{}l}{\small\itshape Table \thetable{} (continued)}\\
\toprule
\textbf{ID} & \textbf{Non-Functional Requirement} & \textbf{Acceptance Criteria} \\
\midrule
\endhead
\midrule
\multicolumn{3}{r@{}}{\footnotesize\itshape Continued on next page}\\
\endfoot
\bottomrule
\endlastfoot
NFR-1.1 & The system shall respond to ordinary user interactions within an acceptable time under normal operating conditions. & TBD \\ \addlinespace[2pt]
NFR-1.2 & The system shall generate and validate short-trip itineraries within an acceptable processing time. & TBD \\ \addlinespace[2pt]
NFR-1.3 & The system shall complete dynamic itinerary recovery within an acceptable processing time. & TBD \\ \addlinespace[2pt]
NFR-1.4 & The system shall provide visible processing feedback for operations that require noticeable waiting time. & Loading or progress indication is displayed during long-running operations. \\ \addlinespace[2pt]
NFR-1.5 & The system shall support the expected number of concurrent users without unacceptable performance degradation. & TBD \\ \addlinespace[2pt]
NFR-1.6 & The system shall avoid unnecessary repeated requests to external mapping and routing services. & Repeated calls are reduced through appropriate request management. \\ \addlinespace[2pt]
\end{longtable}
\endgroup


\Needspace*{12\baselineskip}
\subsection{NFR-2: Usability and Accessibility}

\begingroup
\small
\setlength{\tabcolsep}{5pt}
\renewcommand{\arraystretch}{1.18}
\begin{longtable}{@{}P{\dimexpr 0.15\linewidth-0.60\tabcolsep\relax}P{\dimexpr 0.49\linewidth-1.96\tabcolsep\relax}P{\dimexpr 0.36\linewidth-1.44\tabcolsep\relax}@{}}
\caption{Usability and Accessibility}\label{tab:req-14}\\
\toprule
\textbf{ID} & \textbf{Non-Functional Requirement} & \textbf{Acceptance Criteria} \\
\midrule
\endfirsthead
\multicolumn{3}{@{}l}{\small\itshape Table \thetable{} (continued)}\\
\toprule
\textbf{ID} & \textbf{Non-Functional Requirement} & \textbf{Acceptance Criteria} \\
\midrule
\endhead
\midrule
\multicolumn{3}{r@{}}{\footnotesize\itshape Continued on next page}\\
\endfoot
\bottomrule
\endlastfoot
NFR-2.1 & The system shall provide a consistent and understandable user interface across supported devices. & Core functions remain accessible on supported screen sizes. \\ \addlinespace[2pt]
NFR-2.2 & The system shall minimize unnecessary input during trip creation by using available defaults and reusable preferences. & Previously configured defaults are applied where relevant. \\ \addlinespace[2pt]
NFR-2.3 & The system shall clearly present activity schedules, locations, travel durations, and expected costs. & Required trip information is visible in the Trip Summary. \\ \addlinespace[2pt]
NFR-2.4 & The system shall provide understandable feedback for invalid inputs and infeasible itineraries. & Validation feedback identifies the relevant issue. \\ \addlinespace[2pt]
NFR-2.5 & The system shall provide clear and accessible controls for trip execution actions. & Users can identify and perform supported actions. \\ \addlinespace[2pt]
NFR-2.6 & The system shall distinguish between planned activities, completed activities, and proposed itinerary adjustments. & Different activity states are identifiable in the interface. \\ \addlinespace[2pt]
NFR-2.7 & The system shall support touch-friendly interaction on mobile devices. & Touch controls, map interactions, and forms remain usable on supported devices. \\ \addlinespace[2pt]
NFR-2.8 & The system shall support efficient trip creation for first-time users. & Usability testing target: TBD \\ \addlinespace[2pt]
\end{longtable}
\endgroup


\Needspace*{12\baselineskip}
\subsection{NFR-3: Reliability and Availability}

\begingroup
\small
\setlength{\tabcolsep}{5pt}
\renewcommand{\arraystretch}{1.18}
\begin{longtable}{@{}P{\dimexpr 0.15\linewidth-0.60\tabcolsep\relax}P{\dimexpr 0.49\linewidth-1.96\tabcolsep\relax}P{\dimexpr 0.36\linewidth-1.44\tabcolsep\relax}@{}}
\caption{Reliability and Availability}\label{tab:req-15}\\
\toprule
\textbf{ID} & \textbf{Non-Functional Requirement} & \textbf{Acceptance Criteria} \\
\midrule
\endfirsthead
\multicolumn{3}{@{}l}{\small\itshape Table \thetable{} (continued)}\\
\toprule
\textbf{ID} & \textbf{Non-Functional Requirement} & \textbf{Acceptance Criteria} \\
\midrule
\endhead
\midrule
\multicolumn{3}{r@{}}{\footnotesize\itshape Continued on next page}\\
\endfoot
\bottomrule
\endlastfoot
NFR-3.1 & The system shall maintain an acceptable level of availability during normal operation. & TBD \\ \addlinespace[2pt]
NFR-3.2 & The system shall preserve successfully saved trip data and confirmed progress events. & Saved data remains available after reauthentication or application reload. \\ \addlinespace[2pt]
NFR-3.3 & The system shall prevent inconsistent trip states when itinerary updates fail or are interrupted. & Failed operations do not leave partially committed trip changes. \\ \addlinespace[2pt]
NFR-3.4 & The system shall handle unavailable external services without causing unrelated functionality to fail. & Unaffected functions remain accessible. \\ \addlinespace[2pt]
NFR-3.5 & The system shall clearly report when required external information cannot be retrieved or validated. & An understandable error is displayed. \\ \addlinespace[2pt]
NFR-3.6 & The system shall preserve completed activity history during replanning. & Completed records remain unchanged after an accepted replan. \\ \addlinespace[2pt]
NFR-3.7 & The system shall handle repeated or interrupted trip execution requests consistently. & Repeated requests do not produce duplicate progress transitions. \\ \addlinespace[2pt]
\end{longtable}
\endgroup


\Needspace*{12\baselineskip}
\subsection{NFR-4: Security and Privacy}

\begingroup
\small
\setlength{\tabcolsep}{5pt}
\renewcommand{\arraystretch}{1.18}
\begin{longtable}{@{}P{\dimexpr 0.15\linewidth-0.60\tabcolsep\relax}P{\dimexpr 0.49\linewidth-1.96\tabcolsep\relax}P{\dimexpr 0.36\linewidth-1.44\tabcolsep\relax}@{}}
\caption{Security and Privacy}\label{tab:req-16}\\
\toprule
\textbf{ID} & \textbf{Non-Functional Requirement} & \textbf{Acceptance Criteria} \\
\midrule
\endfirsthead
\multicolumn{3}{@{}l}{\small\itshape Table \thetable{} (continued)}\\
\toprule
\textbf{ID} & \textbf{Non-Functional Requirement} & \textbf{Acceptance Criteria} \\
\midrule
\endhead
\midrule
\multicolumn{3}{r@{}}{\footnotesize\itshape Continued on next page}\\
\endfoot
\bottomrule
\endlastfoot
NFR-4.1 & The system shall protect communication involving authentication and personal information using secure transport protocols. & HTTPS is used in deployment. \\ \addlinespace[2pt]
NFR-4.2 & The system shall securely protect user credentials. & Passwords are stored using an appropriate password-hashing mechanism. \\ \addlinespace[2pt]
NFR-4.3 & The system shall enforce authorization for protected operations. & Unauthorized requests are rejected. \\ \addlinespace[2pt]
NFR-4.4 & The system shall prevent unauthorized access to personal trip records and preferences. & Cross-account access attempts are denied. \\ \addlinespace[2pt]
NFR-4.5 & The system shall validate untrusted input received through user interfaces and external services. & Invalid or malicious input is rejected or safely handled. \\ \addlinespace[2pt]
NFR-4.6 & The system shall protect authentication tokens and external service credentials. & Secrets are not exposed in public client code or application logs. \\ \addlinespace[2pt]
NFR-4.7 & The system shall limit the collection and retention of personal and location information to what is necessary for supported functionality. & Data handling follows the defined privacy policy. \\ \addlinespace[2pt]
NFR-4.8 & The system shall not depend on continuous background location tracking for core trip execution. & Core progress tracking works through explicit user actions. \\ \addlinespace[2pt]
\end{longtable}
\endgroup


\Needspace*{12\baselineskip}
\subsection{NFR-5: Maintainability and Modifiability}

\begingroup
\small
\setlength{\tabcolsep}{5pt}
\renewcommand{\arraystretch}{1.18}
\begin{longtable}{@{}P{\dimexpr 0.15\linewidth-0.60\tabcolsep\relax}P{\dimexpr 0.49\linewidth-1.96\tabcolsep\relax}P{\dimexpr 0.36\linewidth-1.44\tabcolsep\relax}@{}}
\caption{Maintainability and Modifiability}\label{tab:req-17}\\
\toprule
\textbf{ID} & \textbf{Non-Functional Requirement} & \textbf{Acceptance Criteria} \\
\midrule
\endfirsthead
\multicolumn{3}{@{}l}{\small\itshape Table \thetable{} (continued)}\\
\toprule
\textbf{ID} & \textbf{Non-Functional Requirement} & \textbf{Acceptance Criteria} \\
\midrule
\endhead
\midrule
\multicolumn{3}{r@{}}{\footnotesize\itshape Continued on next page}\\
\endfoot
\bottomrule
\endlastfoot
NFR-5.1 & The system shall maintain clear separation between user management, trip planning, feasibility validation, trip execution, and external service integration. & Changes to one functional area do not require unrelated business logic changes. \\ \addlinespace[2pt]
NFR-5.2 & The system shall support modifications to itinerary-generation strategies without requiring major changes to account management or trip storage. & Planning logic can be changed behind stable functional interfaces. \\ \addlinespace[2pt]
NFR-5.3 & The system shall isolate provider-specific mapping and routing integration details from core planning logic. & External service adapters can be replaced without changing core domain rules. \\ \addlinespace[2pt]
NFR-5.4 & The system shall reuse core planning and feasibility rules across itinerary generation, manual editing, and replanning. & These workflows apply consistent validation behavior. \\ \addlinespace[2pt]
NFR-5.5 & The system shall provide sufficient diagnostic information for troubleshooting. & Relevant operational errors are logged. \\ \addlinespace[2pt]
NFR-5.6 & The system shall support automated testing of critical functions. & Authentication, validation, trip lifecycle, and replanning behavior can be tested independently. \\ \addlinespace[2pt]
NFR-5.7 & The mobile interface shall reuse existing application business logic and backend services rather than introduce independent planning rules. & Equivalent desktop and mobile requests produce consistent results. \\ \addlinespace[2pt]
\end{longtable}
\endgroup


\Needspace*{12\baselineskip}
\subsection{NFR-6: Scalability}

\begingroup
\small
\setlength{\tabcolsep}{5pt}
\renewcommand{\arraystretch}{1.18}
\begin{longtable}{@{}P{\dimexpr 0.15\linewidth-0.60\tabcolsep\relax}P{\dimexpr 0.49\linewidth-1.96\tabcolsep\relax}P{\dimexpr 0.36\linewidth-1.44\tabcolsep\relax}@{}}
\caption{Scalability}\label{tab:req-18}\\
\toprule
\textbf{ID} & \textbf{Non-Functional Requirement} & \textbf{Acceptance Criteria} \\
\midrule
\endfirsthead
\multicolumn{3}{@{}l}{\small\itshape Table \thetable{} (continued)}\\
\toprule
\textbf{ID} & \textbf{Non-Functional Requirement} & \textbf{Acceptance Criteria} \\
\midrule
\endhead
\midrule
\multicolumn{3}{r@{}}{\footnotesize\itshape Continued on next page}\\
\endfoot
\bottomrule
\endlastfoot
NFR-6.1 & The system shall support growth in registered users and stored itineraries without requiring a complete architectural redesign. & Target workload: TBD \\ \addlinespace[2pt]
NFR-6.2 & The system shall support increased itinerary-generation workloads through appropriate resource management. & Target capacity: TBD \\ \addlinespace[2pt]
NFR-6.3 & The system shall manage external API consumption according to provider limitations. & Provider rate limits are respected. \\ \addlinespace[2pt]
NFR-6.4 & The system shall avoid unnecessary duplication of persistent trip data and planning computations. & Shared trip representations and planning services are reused. \\ \addlinespace[2pt]
\end{longtable}
\endgroup


\Needspace*{12\baselineskip}
\subsection{NFR-7: Interoperability}

\begingroup
\small
\setlength{\tabcolsep}{5pt}
\renewcommand{\arraystretch}{1.18}
\begin{longtable}{@{}P{\dimexpr 0.15\linewidth-0.60\tabcolsep\relax}P{\dimexpr 0.49\linewidth-1.96\tabcolsep\relax}P{\dimexpr 0.36\linewidth-1.44\tabcolsep\relax}@{}}
\caption{Interoperability}\label{tab:req-19}\\
\toprule
\textbf{ID} & \textbf{Non-Functional Requirement} & \textbf{Acceptance Criteria} \\
\midrule
\endfirsthead
\multicolumn{3}{@{}l}{\small\itshape Table \thetable{} (continued)}\\
\toprule
\textbf{ID} & \textbf{Non-Functional Requirement} & \textbf{Acceptance Criteria} \\
\midrule
\endhead
\midrule
\multicolumn{3}{r@{}}{\footnotesize\itshape Continued on next page}\\
\endfoot
\bottomrule
\endlastfoot
NFR-7.1 & The system shall integrate with external geographic services for place search and route information. & Supported external service requests and responses are handled correctly. \\ \addlinespace[2pt]
NFR-7.2 & The system shall use defined data contracts for communication between the frontend, backend, and optional AI service. & Service responses conform to documented schemas. \\ \addlinespace[2pt]
NFR-7.3 & The system shall normalize external place and route data before using it in core planning operations. & Provider-specific formats are converted into internal representations. \\ \addlinespace[2pt]
NFR-7.4 & The system shall handle incomplete or unavailable external service responses. & Invalid external responses do not produce silently accepted infeasible plans. \\ \addlinespace[2pt]
NFR-7.5 & The system shall support opening routes or destinations in an external navigation application. & Navigation handoff is available for supported destinations. \\ \addlinespace[2pt]
\end{longtable}
\endgroup


\Needspace*{12\baselineskip}
\subsection{NFR-8: Portability and Responsive Compatibility}

\begingroup
\small
\setlength{\tabcolsep}{5pt}
\renewcommand{\arraystretch}{1.18}
\begin{longtable}{@{}P{\dimexpr 0.15\linewidth-0.60\tabcolsep\relax}P{\dimexpr 0.49\linewidth-1.96\tabcolsep\relax}P{\dimexpr 0.36\linewidth-1.44\tabcolsep\relax}@{}}
\caption{Portability and Responsive Compatibility}\label{tab:req-20}\\
\toprule
\textbf{ID} & \textbf{Non-Functional Requirement} & \textbf{Acceptance Criteria} \\
\midrule
\endfirsthead
\multicolumn{3}{@{}l}{\small\itshape Table \thetable{} (continued)}\\
\toprule
\textbf{ID} & \textbf{Non-Functional Requirement} & \textbf{Acceptance Criteria} \\
\midrule
\endhead
\midrule
\multicolumn{3}{r@{}}{\footnotesize\itshape Continued on next page}\\
\endfoot
\bottomrule
\endlastfoot
NFR-8.1 & The system shall provide a responsive interface for supported desktop and mobile screen sizes. & Core workflows remain usable at representative viewport sizes. \\ \addlinespace[2pt]
NFR-8.2 & The mobile interface shall prioritize map visibility during itinerary viewing and active trip execution. & Secondary content does not permanently obstruct the primary map workspace. \\ \addlinespace[2pt]
NFR-8.3 & The mobile interface shall provide access to the full Trip Summary without requiring a separate business workflow. & Users can open and close the summary while retaining the current trip context. \\ \addlinespace[2pt]
NFR-8.4 & The system shall support touch-based place selection, map gestures, and itinerary editing. & Core map and editing interactions work on supported touch devices. \\ \addlinespace[2pt]
NFR-8.5 & The system shall preserve consistent trip data and behavior across desktop and mobile interfaces. & The same saved trip produces equivalent authoritative information on both interfaces. \\ \addlinespace[2pt]
NFR-8.6 & The mobile interface shall remain usable with on-screen keyboards, scrolling, overlays, and long itineraries. & Device interaction testing: TBD \\ \addlinespace[2pt]
\end{longtable}
\endgroup


\Needspace*{12\baselineskip}
\subsection{NFR-9: Data Integrity and Consistency}

\begingroup
\small
\setlength{\tabcolsep}{5pt}
\renewcommand{\arraystretch}{1.18}
\begin{longtable}{@{}P{\dimexpr 0.15\linewidth-0.60\tabcolsep\relax}P{\dimexpr 0.49\linewidth-1.96\tabcolsep\relax}P{\dimexpr 0.36\linewidth-1.44\tabcolsep\relax}@{}}
\caption{Data Integrity and Consistency}\label{tab:req-21}\\
\toprule
\textbf{ID} & \textbf{Non-Functional Requirement} & \textbf{Acceptance Criteria} \\
\midrule
\endfirsthead
\multicolumn{3}{@{}l}{\small\itshape Table \thetable{} (continued)}\\
\toprule
\textbf{ID} & \textbf{Non-Functional Requirement} & \textbf{Acceptance Criteria} \\
\midrule
\endhead
\midrule
\multicolumn{3}{r@{}}{\footnotesize\itshape Continued on next page}\\
\endfoot
\bottomrule
\endlastfoot
NFR-9.1 & The system shall maintain consistent trip and itinerary data across supported user interfaces. & Desktop and mobile retrieve consistent server-side trip records. \\ \addlinespace[2pt]
NFR-9.2 & The system shall ensure that confirmed itinerary modifications are validated before becoming authoritative. & Invalid changes are not committed as accepted plans. \\ \addlinespace[2pt]
NFR-9.3 & The system shall preserve completed activity records when the unfinished itinerary is modified. & Completed activity history remains unchanged. \\ \addlinespace[2pt]
NFR-9.4 & The system shall maintain valid trip lifecycle transitions. & Unsupported state transitions are rejected. \\ \addlinespace[2pt]
NFR-9.5 & The system shall distinguish temporary guest data from persistently saved user data. & Guest data is not treated as permanently saved until the required authentication and save process succeeds. \\ \addlinespace[2pt]
\end{longtable}
\endgroup


\Needspace*{12\baselineskip}
\subsection{NFR-10: AI Service Reliability and Output Validation}

\begingroup
\small
\setlength{\tabcolsep}{5pt}
\renewcommand{\arraystretch}{1.18}
\begin{longtable}{@{}P{\dimexpr 0.15\linewidth-0.60\tabcolsep\relax}P{\dimexpr 0.49\linewidth-1.96\tabcolsep\relax}P{\dimexpr 0.36\linewidth-1.44\tabcolsep\relax}@{}}
\caption{AI Service Reliability and Output Validation}\label{tab:req-22}\\
\toprule
\textbf{ID} & \textbf{Non-Functional Requirement} & \textbf{Acceptance Criteria} \\
\midrule
\endfirsthead
\multicolumn{3}{@{}l}{\small\itshape Table \thetable{} (continued)}\\
\toprule
\textbf{ID} & \textbf{Non-Functional Requirement} & \textbf{Acceptance Criteria} \\
\midrule
\endhead
\midrule
\multicolumn{3}{r@{}}{\footnotesize\itshape Continued on next page}\\
\endfoot
\bottomrule
\endlastfoot
NFR-10.1 & The system shall treat AI-generated outputs as untrusted suggestions until they have been structurally validated. & Invalid structured responses are rejected. \\ \addlinespace[2pt]
NFR-10.2 & The system shall validate AI-proposed itinerary changes using the core planning and feasibility functionality. & AI suggestions cannot bypass mandatory constraints. \\ \addlinespace[2pt]
NFR-10.3 & The system shall handle AI service failures without disabling the core non-AI planning workflow. & Non-AI trip creation and editing remain available. \\ \addlinespace[2pt]
NFR-10.4 & The system shall provide understandable feedback when AI processing fails or produces unsupported instructions. & Users receive a clear failure or correction message. \\ \addlinespace[2pt]
NFR-10.5 & The system shall protect sensitive information exchanged with external AI providers according to the defined privacy policy. & Data-sharing and retention rules: TBD \\ \addlinespace[2pt]
NFR-10.6 & The system shall process supported AI requests within an acceptable response time. & TBD \\ \addlinespace[2pt]
\end{longtable}
\endgroup

\clearpage
\section{System Models}
\label{sec:system-models}

\subsection{Use Case Table}
\label{sec:use-case-table}
\textbf{TBD.}

\subsection{Use Case Diagram}
\label{sec:use-case-diagram}
\textbf{TBD.}

\subsection{Activity Diagram}
\label{sec:activity-diagram}
\textbf{TBD.}

\subsection{Sequence Diagram}
\label{sec:sequence-diagram}
\textbf{TBD.}

\subsection{Class Diagram}
\label{sec:class-diagram}
\textbf{TBD.}
